@ -0,0 +1,317 @@
%%%
%
% $Autor: Shubhankar Dumka, Gautam Ramesh, Siddhanth Sharma  $
% $Datum: 2025-12-01 $
% $Dateiname: 
% $Version: 4620 $
%
% !TeX spellcheck = GB
% !TeX program = pdflatex
% !TeX encoding = utf8
%
%%%

\chapter{\textcolor{red}{Methodology}}

\index{Methodology}

% --------------------------------------------------------
\section{Introduction}
\index{Methodology!Introduction}

\noindent
Convolutional Neural Networks (CNNs) are a foundational class of deep learning models specifically designed to process data with a grid-like topology, such as images or sequential signals. At their core, CNNs utilize the mathematical operation of convolution to extract local spatial features by sliding learnable kernels across the input, enabling the model to detect elementary patterns such as edges, corners, and textures \cite{goodfellow2016deep}. These convolutional layers are often complemented by pooling operations, which downsample feature maps to reduce computational complexity while providing a degree of translational invariance \cite{insights2018cnn}. By integrating these hierarchical operations, CNNs form increasingly abstract representations of the input data, allowing the network to learn robust features directly from raw pixel intensities without the need for handcrafted descriptors \cite{purwono2020understanding}.


% --------------------------------------------------------
\section{CRISP-DM}
\index{Methodology!CRISP-DM}

\section{Methodology: CRISP-DM for License Plate Detection on Jetson Nano}

\noindent
This project follows the \textbf{CRISP-DM (Cross-Industry Standard Process for Data Mining)} methodology, a widely adopted and structured framework for data-driven system development \cite{chapman2000crisp}. CRISP-DM provides a systematic approach consisting of six iterative phases: business understanding, data understanding, data preparation, modeling, evaluation, and deployment. Its flexibility and domain-independence make it particularly suitable for embedded computer vision applications, where machine learning models must be trained, evaluated, and deployed under hardware constraints \cite{shearer2000crispdm}. The following subsections describe how each CRISP-DM phase was applied in the context of license plate detection and timestamp-based fee calculation using the NVIDIA Jetson Nano.

\subsection{Business Understanding}

The primary objective of this project is to design and implement a real-time license plate detection system capable of operating on embedded hardware. The system must detect vehicle license plates, extract alphanumeric characters, record corresponding entry and exit timestamps, and store this information in a SQLite database for subsequent fee calculation based on parking duration. Such systems are commonly used in intelligent transportation systems, automated parking management, and access control applications \cite{anagnostopoulos2008license,du2013automatic}. 

From a business perspective, the key requirements include reliable detection accuracy, real-time performance, low operational cost, and independence from cloud-based processing. The NVIDIA Jetson Nano was selected as the target platform due to its CUDA-enabled GPU, low power consumption, and suitability for edge AI applications \cite{nvidia_jetson}. These constraints guided the selection of lightweight detection models, efficient image-processing pipelines, and embedded-friendly database solutions.

\subsection{Data Understanding}

The dataset used in this project consists of 1,243 labeled images containing German vehicle license plates captured under varying lighting conditions, viewing angles, and background complexity. Understanding the characteristics of the dataset was critical to selecting appropriate preprocessing steps and model architectures. Variations in plate size, font style, illumination, and occlusion are well-known challenges in automatic license plate recognition (ALPR) systems \cite{silva2018license}. 

Initial exploratory analysis focused on image resolution, aspect ratios, and class distribution to ensure suitability for training an object detection model. Sample images were visually inspected using OpenCV to identify noise, blur, and contrast issues that could negatively affect detection and recognition performance \cite{bradski2000opencv}.

\subsection{Data Preparation}

During the data preparation phase, images were loaded and processed using the OpenCV (\texttt{cv2}) library. Images were resized to match the input requirements of the YOLO-based detection model \cite{opencv_wiki}. 

Data augmentation techniques such as scaling, translation, and brightness adjustment were applied to improve model robustness and generalization, particularly given the relatively limited dataset size \cite{shorten2019survey}. These steps are essential in computer vision tasks to mitigate overfitting and improve performance under real-world conditions, especially in embedded deployment scenarios.

\subsection{Modeling}

For the modeling phase, a YOLO-based object detection architecture was selected due to its balance between detection accuracy and real-time inference speed \cite{redmon2016you,bochkovskiy2020yolov4}. YOLO models perform object localization and classification in a single forward pass, making them well suited for edge devices such as the Jetson Nano. The model was trained to detect license plate bounding boxes using the prepared dataset.

Following license plate localization, character recognition was performed using EasyOCR, an open-source optical character recognition library designed to work effectively with natural scene text \cite{jaidedai2020easyocr}. The integration of YOLO for detection and EasyOCR for character extraction enabled a modular pipeline where detection and recognition could be optimized independently. Training was conducted on a workstation with GPU acceleration, after which the trained model weights were transferred to the Jetson Nano for inference.

\subsection{Evaluation}

Model evaluation focused on both detection accuracy and system-level performance. Detection quality was assessed using standard object detection metrics such as precision and recall, while OCR output was evaluated by comparing detected characters against ground truth labels \cite{everingham2010pascal}. In addition to accuracy, inference latency and frame rate were measured on the Jetson Nano to ensure that real-time constraints were satisfied.

End-to-end testing was performed by capturing live video streams from an external camera connected to the Jetson Nano, detecting license plates, extracting text, and recording timestamps. This holistic evaluation approach ensured that the integrated system met both functional and performance requirements under realistic operating conditions.

\subsection{Deployment}

In the deployment phase, the trained YOLO model and OCR pipeline were integrated into a real-time application running on the Jetson Nano. OpenCV was used for camera interfacing and frame preprocessing, while PyTorch handled model inference on the CUDA-enabled GPU \cite{paszke2019pytorch}. Detected license plate numbers, along with entry and exit timestamps, were stored in a lightweight SQLite database, chosen for its simplicity, reliability, and suitability for embedded systems \cite{sqlite_doc}.

The deployed system continuously monitors incoming frames, updates database records, and enables subsequent computation of parking duration and fee calculation. This deployment strategy demonstrates the feasibility of performing end-to-end license plate detection, recognition, and data management entirely on edge hardware, reducing latency and eliminating dependence on cloud infrastructure \cite{shi2016edge}.


\section{KDD Process}
\index{Methodology!KDD}

\section{Methodology: Knowledge Discovery in Databases (KDD)}

\noindent
The Knowledge Discovery in Databases (KDD) process provides a structured framework for transforming raw data into meaningful and actionable knowledge through a sequence of well-defined stages \cite{fayyad1996knowledge}. In this project, KDD is applied holistically across both the \textit{model training phase} and the \textit{operational deployment phase}. Unlike methodologies that focus solely on post-deployment analytics, this consolidated KDD approach captures how knowledge is progressively extracted from raw visual data, beginning with dataset construction and model learning, and culminating in real-time license plate analysis and fee computation on embedded hardware.

\begin{figure}[htbp]
	\centering
	\begin{tikzpicture}[
		node distance=1.6cm, % increased vertical spacing
		every node/.style={
			rectangle,
			draw=black,
			thick,
			align=center,
			minimum width=3.0cm,
			minimum height=0.85cm,
			font=\scriptsize,
			fill=gray!10
		},
		header/.style={
			font=\footnotesize\bfseries,
			draw=none,
			fill=none
		},
		arrow/.style={
			->,
			thick
		}
		]
		
		% Column headers
		\node[header] (train_h) {Training Phase};
		\node[header, right=3.8cm of train_h] (deploy_h) {Deployment Phase};
		
		% KDD-1
		\node (s1) [below=1.0cm of train_h] {KDD--1\\Data Selection\\Labeled Images};
		\node (s2) [below=1.0cm of deploy_h] {KDD--1\\Data Selection\\Live Camera Frames};
		
		% KDD-2
		\node (p1) [below of=s1] {KDD--2\\Preprocessing\\Resize, Normalize, Augment};
		\node (p2) [below of=s2] {KDD--2\\Preprocessing\\Frame Conditioning};
		
		% KDD-3
		\node (t1) [below of=p1] {KDD--3\\Transformation\\Image--Label Pairs};
		\node (t2) [below of=p2] {KDD--3\\Transformation\\Bounding Boxes, OCR};
		
		% KDD-4
		\node (m1) [below of=t1] {KDD--4\\Data Mining\\YOLO Model Training};
		\node (m2) [below of=t2] {KDD--4\\Data Mining\\Real-Time Inference};
		
		% KDD-5
		\node (e1) [below of=m1] {KDD--5\\Evaluation\\Accuracy, Loss Metrics};
		\node (e2) [below of=m2] {KDD--5\\Evaluation\\Latency, OCR Quality};
		
		% KDD-6
		\node (k1) [below of=e1] {KDD--6\\Knowledge Storage\\Trained Model};
		\node (k2) [below of=e2] {KDD--6\\Knowledge Storage\\SQLite Records};
		
		% Training arrows (downward)
		\draw[arrow] (s1.south) -- (p1.north);
		\draw[arrow] (p1.south) -- (t1.north);
		\draw[arrow] (t1.south) -- (m1.north);
		\draw[arrow] (m1.south) -- (e1.north);
		\draw[arrow] (e1.south) -- (k1.north);
		
		% Deployment arrows (downward)
		\draw[arrow] (s2.south) -- (p2.north);
		\draw[arrow] (p2.south) -- (t2.north);
		\draw[arrow] (t2.south) -- (m2.north);
		\draw[arrow] (m2.south) -- (e2.north);
		\draw[arrow] (e2.south) -- (k2.north);
		
		% Legend
		\node[draw, thick, align=left, font=\scriptsize, fill=white,
		below=1.8cm of k1, minimum width=7.2cm] (legend)
		{\textbf{Legend:}\\
			KDD--1: Data Selection \\
			KDD--2: Preprocessing and Cleaning \\
			KDD--3: Data Transformation \\
			KDD--4: Data Mining \\
			KDD--5: Interpretation and Evaluation \\
			KDD--6: Knowledge Presentation and Storage};
		
	\end{tikzpicture}
	\caption{Vertical representation of the Knowledge Discovery in Databases (KDD) process applied to both training and deployment stages of the license plate detection system. Arrows indicate downward progression through each KDD phase.}
	\label{fig:kdd_vertical_refined}
\end{figure}
 


\subsection{Data Selection}

The KDD process begins with data selection, where relevant data sources are identified and collected. For the training phase, a dataset consisting of 1,243 labeled images containing German vehicle license plates was selected. These images represent diverse environmental conditions, including variations in illumination, plate orientation, distance, and background complexity. Such diversity is critical for learning robust visual patterns in automatic license plate recognition (ALPR) systems \cite{du2013automatic,silva2018license}.

During deployment, data selection occurs dynamically, as image frames are captured in real time from an external camera connected to the NVIDIA Jetson Nano. Only frames containing vehicles or candidate license plates are forwarded to subsequent stages, ensuring computational efficiency and relevance of the extracted data.

\subsection{Data Preprocessing and Cleaning}

In both training and deployment stages, preprocessing and cleaning play a central role in improving data quality. During training, images are loaded using the OpenCV library and undergo resizing, normalization, and color space conversion to conform to the input requirements of the YOLO-based detection model \cite{bradski2000opencv,opencv_wiki}. Noisy or low-quality samples are either corrected or excluded to prevent degradation of the learning process.

To enhance generalization and mitigate overfitting, data augmentation techniques such as geometric transformations and brightness variation are applied during training \cite{shorten2019survey}. During deployment, preprocessing ensures consistency between live camera input and the data distribution observed during training, which is essential for maintaining reliable inference performance on embedded platforms.

\subsection{Data Transformation}

The transformation stage converts preprocessed data into representations suitable for knowledge extraction. During training, annotated images are transformed into structured input-output pairs consisting of pixel data and corresponding bounding box labels. These representations enable supervised learning of spatial and semantic patterns associated with license plates \cite{redmon2016you}.

In the deployment phase, transformation occurs through inference. Raw image frames are converted into bounding box predictions by the trained YOLO model, effectively transforming unstructured pixel data into structured object-level information \cite{bochkovskiy2020yolov4}. Detected license plate regions are further transformed into alphanumeric text using optical character recognition, bridging the gap between visual perception and symbolic data representation \cite{jaidedai2020easyocr}.

\subsection{Data Mining}

Data mining represents the core analytical stage of the KDD process, where algorithms extract patterns or models from transformed data \cite{han2011data}. In the training phase, data mining is realized through deep learning model optimization. The YOLO detection model learns hierarchical visual features that distinguish license plates from background elements by iteratively minimizing a loss function over the training dataset.

This learning process embodies pattern discovery at scale, as the model encodes complex spatial and contextual relationships within its parameters \cite{lecun2015deep}. In the deployment phase, the trained models continue to perform data mining in an online manner, applying learned patterns to incoming image streams to detect license plates and recognize characters in real time.

\subsection{Interpretation and Evaluation}

Interpretation and evaluation occur at multiple stages of the KDD pipeline. During training, model performance is evaluated using standard object detection and recognition metrics, such as precision, recall, and qualitative inspection of predicted bounding boxes \cite{everingham2010pascal}. These evaluations guide hyperparameter tuning and model selection, ensuring that the discovered patterns generalize beyond the training dataset.

During deployment, interpretation focuses on the correctness and consistency of extracted license plate numbers and timestamps. Detection errors, OCR inaccuracies, and duplicate records are analyzed and filtered where necessary to maintain data integrity. System-level performance metrics, including inference latency and reliability on the Jetson Nano, are also evaluated to ensure operational feasibility.

\subsection{Knowledge Presentation and Storage}

The final stage of the KDD process involves presenting and storing extracted knowledge in a form suitable for practical use. In the deployment phase, recognized license plate numbers are paired with entry and exit timestamps and stored in a SQLite database. SQLite is well suited for this task due to its lightweight architecture, transactional support, and compatibility with embedded systems \cite{sqlite_doc}.

Higher-level knowledge is subsequently derived by computing temporal differences between entry and exit events and applying predefined pricing rules to calculate parking fees. This transformation of low-level visual observations into actionable financial information represents the culmination of the KDD process, enabling autonomous decision-making at the edge without reliance on cloud-based infrastructure \cite{shi2016edge}.

\noindent
By consolidating both the training and operational stages within the KDD framework, this methodology demonstrates how knowledge discovery extends beyond post-deployment analytics and encompasses the entire lifecycle of a data-driven computer vision system.


% --------------------------------------------------------
\section{Machine Learning Pipeline}
\index{Methodology!MLPipeline}

\noindent The machine learning pipeline in this project is designed to enable real-time license plate detection and recognition on the NVIDIA Jetson Nano. This pipeline follows a structured sequence of stages inspired by conventional machine learning workflows \cite{kuhn2013applied,goodfellow2016deep} and is adapted to the constraints of embedded deployment. The pipeline integrates data acquisition, preprocessing, model training, evaluation, deployment, and continuous inference to ensure robust and efficient system operation.

\subsection{Data Acquisition and Understanding}
\noindent The first stage of the pipeline involves collecting and understanding the dataset. For training, 1,243 labeled images of German license plates were used, encompassing variations in illumination, plate orientation, occlusion, and background complexity \cite{du2013automatic,silva2018license}.

Exploratory data analysis (EDA) was performed to:

\begin{itemize}
	\item Inspect image resolution and aspect ratios to ensure compatibility with the YOLO input dimensions \cite{redmon2016you}.
	\item Examine class balance and diversity of license plate appearances to mitigate biases during model training.
	\item Identify noise, blur, or low-contrast images using OpenCV visualization tools \cite{bradski2000opencv}.
\end{itemize}

This stage ensured a comprehensive understanding of the dataset and informed subsequent preprocessing and augmentation decisions.

\subsection{Data Preprocessing}
\noindent Preprocessing was performed to standardize input data and improve model performance \cite{shorten2019survey}. Key steps included:

\begin{itemize}
	\item Image resizing and normalization to match the YOLO model input requirements.
	\item Data augmentation using techniques such as scaling, translation, rotation, and brightness adjustment to enhance robustness against real-world variability.
	\item Color space conversion (RGB to grayscale where appropriate) to optimize OCR performance with EasyOCR \cite{jaidedai2020easyocr}.
\end{itemize}

During deployment, incoming frames from the external camera were preprocessed in real time to match the data distribution observed during training, ensuring consistent inference performance on the Jetson Nano.

\subsection{Feature Extraction and Representation}
\noindent Unlike traditional machine learning pipelines relying on handcrafted features, this pipeline leverages Convolutional Neural Networks (CNNs) within the YOLO architecture to automatically learn hierarchical visual features from raw image pixels \cite{goodfellow2016deep,purwono2020understanding}.

The YOLO model encodes spatial and semantic information via convolutional layers, producing bounding box predictions for license plates in a single forward pass. Subsequently, extracted regions of interest (ROIs) are passed to EasyOCR, which converts visual patterns into alphanumeric characters, effectively transforming raw pixel data into structured symbolic representations \cite{jaidedai2020easyocr}.

\subsection{Model Training and Optimization}
\noindent The YOLO model was trained on the preprocessed dataset using GPU-accelerated workstations before deployment on the Jetson Nano. Training involved:

\begin{itemize}
	\item Loss minimization over bounding box coordinates and object confidence scores, following standard YOLO objectives \cite{bochkovskiy2020yolov4}.
	\item Hyperparameter tuning for learning rate, batch size, and data augmentation parameters to maximize detection accuracy while preserving computational efficiency \cite{bergstra2012random}.
	\item Regularization techniques such as dropout and early stopping to mitigate overfitting given the limited dataset size.
\end{itemize}

This stage ensured that the trained model could generalize effectively to unseen license plates under diverse environmental conditions.

\subsection{Model Evaluation}
\noindent Model performance was assessed through multiple metrics:

\begin{itemize}
	\item Detection metrics: Precision, recall, and mean Average Precision (mAP) for bounding box localization \cite{everingham2010pascal}.
	\item Recognition metrics: Character-level accuracy for OCR outputs.
	\item System-level performance: Inference latency, frame rate, and memory usage on the Jetson Nano to ensure real-time processing capabilities.
\end{itemize}

Evaluation was performed iteratively, with errors analyzed to guide further data augmentation, hyperparameter adjustment, and model refinement.

\subsection{Deployment and Inference}
\noindent The trained YOLO and EasyOCR models were deployed on the Jetson Nano, integrating real-time video capture via OpenCV and inference through PyTorch \cite{paszke2019pytorch,bradski2000opencv}. The deployment pipeline consists of:

\begin{enumerate}
	\item Capturing frames from the external camera.
	\item Performing YOLO-based license plate detection.
	\item Applying EasyOCR to detected ROIs for character recognition.
	\item Recording license plate numbers along with entry and exit timestamps in a SQLite database \cite{sqlite_doc}.
	\item Computing parking duration and fees using timestamp differences for subsequent management and billing operations.
\end{enumerate}

This deployment strategy ensures edge-based, low-latency operation without reliance on cloud infrastructure, making the system suitable for embedded parking management applications \cite{shi2016edge}.

\subsection{Continuous Monitoring and Maintenance}
\noindent To maintain accuracy over time, the system supports periodic model evaluation and retraining. New images captured in real-world conditions can be added to the dataset, and the YOLO model can be fine-tuned to account for environmental changes or new license plate styles, following best practices in machine learning lifecycle management \cite{sculley2015hidden}.


% --------------------------------------------------------
\section{Others}
\index{Methodology!Others}

\subsection{Hyperparameters in Convolutional Neural Networks}

Convolutional Neural Networks (CNNs) rely on a variety of hyperparameters that control their learning behavior, feature extraction capability, and generalization performance \cite{goodfellow2016deep, pinto2022springer}. Choosing the right hyperparameters is critical to achieving high accuracy while avoiding overfitting or underfitting. Key hyperparameters include the number of convolutional filters, kernel size, stride, padding, learning rate, batch size, number of epochs, and regularization parameters.

\paragraph{Number of Convolutional Filters:}  
The number of filters \(N_l\) in layer \(l\) determines the number of feature maps generated by the convolution. Each filter \(K_i^{(l)} \in \mathbb{R}^{k \times k}\) detects a specific feature in the input. Mathematically, the convolution operation for filter \(i\) is given by:
\[
F_i^{(l)} = f\big(X^{(l-1)} * K_i^{(l)} + b_i^{(l)}\big)
\]
where \(X^{(l-1)}\) is the input to layer \(l\), \(b_i^{(l)}\) is the bias term, and \(f(\cdot)\) is the activation function. Increasing the number of filters allows the network to capture more complex patterns, but also increases computational cost \cite{pytorchdocs}.

\begin{figure}[H]
	\centering
	\begin{tikzpicture}[block/.style={draw, rectangle, minimum width=2cm, minimum height=1cm, align=center, fill=blue!20}]
		\node[block] (input) {Input Feature Map};
		\node[block, right=1.5cm of input] (conv1) {Filter 1};
		\node[block, above=0.5cm of conv1] (conv2) {Filter 2};
		\node[block, below=0.5cm of conv1] (conv3) {Filter 3};
		\draw[->, thick] (input.east) -- (conv1.west);
		\draw[->, thick] (input.east) -- (conv2.west);
		\draw[->, thick] (input.east) -- (conv3.west);
	\end{tikzpicture}
	\caption{Illustration of multiple convolutional filters generating feature maps from a single input.}
\end{figure}

\paragraph{Kernel Size, Stride, and Padding:}  
The kernel size \(k \times k\) defines the spatial extent of each convolutional filter. Stride \(s\) determines how the kernel moves over the input:
\[
F^{(l)}_{i,j} = f\Big(\sum_{u=0}^{k-1} \sum_{v=0}^{k-1} X^{(l-1)}_{i\cdot s + u, j \cdot s + v} \cdot K^{(l)}_{u,v} + b^{(l)}\Big)
\]
Padding \(p\) adds extra pixels around the input, controlling the output dimensions:
\[
H_{\text{out}} = \frac{H_{\text{in}} - k + 2p}{s} + 1, \quad W_{\text{out}} = \frac{W_{\text{in}} - k + 2p}{s} + 1
\]
Choosing appropriate kernel size, stride, and padding affects feature resolution and the size of the output feature maps \cite{goodfellow2016deep}.

\begin{figure}[H]
	\centering
	\begin{tikzpicture}[block/.style={draw, rectangle, minimum width=2cm, minimum height=1cm, align=center, fill=green!20}]
		\node[block] (input) {Input Image};
		\node[block, right=2cm of input] (kernel) {Kernel $k \times k$};
		\node[block, right=2cm of kernel] (output) {Feature Map};
		\draw[->, thick] (input.east) -- (kernel.west) node[midway, above] {Stride $s$};
		\draw[->, thick] (kernel.east) -- (output.west) node[midway, above] {Padding $p$};
	\end{tikzpicture}
	\caption{Illustration of convolution operation with kernel, stride, and padding.}
\end{figure}

\paragraph{Learning Rate and Optimization:}  
The learning rate \(\eta\) controls the size of the update step in gradient descent:
\[
\theta \leftarrow \theta - \eta \frac{\partial \mathcal{L}}{\partial \theta}
\]
where \(\theta\) represents the network parameters and \(\mathcal{L}\) the loss function. A high learning rate may cause divergence, while a very low learning rate slows training. Optimizers such as SGD, Adam, and RMSProp adjust parameter updates adaptively to improve convergence \cite{kingma2014adam, pytorchdocs}.

\paragraph{Batch Size and Epochs:}  
Batch size \(B\) determines the number of samples processed before updating network weights. The number of epochs \(E\) is the number of complete passes through the training dataset. Smaller batch sizes increase gradient variance but reduce memory usage, while larger batches provide smoother updates. The output update per batch is:
\[
\theta \leftarrow \theta - \eta \frac{1}{B} \sum_{i=1}^{B} \frac{\partial \mathcal{L}(x_i, y_i)}{\partial \theta}
\]
Proper selection of batch size and epochs is crucial to balance convergence speed and generalization \cite{pinto2022springer}.

\paragraph{Regularization Parameters:}  
Techniques like weight decay (L2 regularization) and dropout reduce overfitting. Weight decay adds a penalty term to the loss function:
\[
\mathcal{L}_{\text{total}} = \mathcal{L} + \lambda \sum_i \theta_i^2
\]
where \(\lambda\) is the regularization coefficient. Dropout randomly deactivates a fraction \(p\) of neurons during training, forcing the network to learn more robust features \cite{srivastava2014dropout}.

% --------------------------------------------------------
\section{Methodology Justification}
\index{Methodology!Justification}

\subsection{Requirements for a Convolutional Neural Network}

Convolutional Neural Networks (CNNs) have specific requirements to function efficiently and accurately. One of the primary requirements is \textbf{high-quality labeled data} for supervised learning. The network relies on large datasets of input images paired with ground truth labels to learn meaningful feature representations \cite{goodfellow2016deep}. Insufficient or noisy data can lead to poor generalization and overfitting, making dataset quality and diversity crucial.

Another requirement is \textbf{computational resources}, particularly GPUs, which are essential for accelerating the convolutional and matrix operations involved in training CNNs \cite{pinto2022springer, pytorchdocs}. While CPU-only training is possible, it is significantly slower and may not be practical for large datasets or deep networks. Memory capacity also plays an important role, as intermediate feature maps and batch processing consume significant amounts of RAM.

CNNs also require \textbf{appropriate model architecture and hyperparameters}, such as the number of layers, filter sizes, strides, padding, learning rate, and regularization parameters. Choosing suitable hyperparameters ensures efficient learning, prevents overfitting, and helps the network converge faster during training \cite{krizhevsky2012imagenet, srivastava2014dropout}.

\textbf{Preprocessing and data augmentation} are additional requirements for effective CNN training. Input images often need resizing, normalization, and augmentation techniques such as rotation, flipping, and cropping to improve model robustness \cite{pinto2022springer}. Proper preprocessing ensures the network can handle variations in real-world data, including different lighting conditions, orientations, and background clutter.

Finally, \textbf{Python implementation and software dependencies} are critical for building and deploying CNNs. Frameworks such as PyTorch and TensorFlow provide modules for defining convolutional layers, loss functions, activation functions, and optimizers \cite{pytorchdocs}. Common Python dependencies include:
\begin{itemize}
	\item \texttt{torch} and \texttt{torchvision} for model building and image utilities.
	\item \texttt{numpy} and \texttt{scipy} for numerical operations.
	\item \texttt{matplotlib} or \texttt{seaborn} for visualization.
	\item \texttt{opencv-python} for image capture and preprocessing.
\end{itemize}

For hardware deployment on embedded platforms such as the Jetson Nano, additional requirements include \textbf{CUDA-enabled GPU drivers} and optimized libraries to leverage GPU acceleration \cite{jetsonpytorch}. These software and hardware dependencies enable real-time inference and efficient model training on edge devices, making CNNs practical for applications such as license plate detection.